\begin{scanpage}{0102}{85}
In accordance with (12.17), the effective scattering area of the cylinder is determined with the E-polarization of the incident wave by the function
\[\sigma_E=\pi a^2|\overline\Sigma(\vartheta)|^2,\tag{15.08}\]
and with the H-polarization of the incident wave by the function
\[\sigma_H=\pi a^2|\Sigma(\vartheta)|^2.\tag{15.09}\]
Let us note that Expressions (15.02) for the scattering field may be obtained directly on the basis of an analogy with the Equations (12.06), omitting the calculation of the fields from the uniform and nonuniform parts of the current. In the same way, one may obtain the expressions
\[\begin{aligned}
\overline\Sigma(\vartheta,\vartheta_0)&=[\overline M J_1(\xi)+i\overline N J_2(\xi)]e^{i\frac{kl}2(\cos\vartheta+\cos\vartheta_0)}\\
&\quad-f(3)[J_1(\xi)-iJ_2(\xi)]e^{-i\frac{kl}2(\cos\vartheta+\cos\vartheta_0)},\\
\Sigma(\vartheta,\vartheta_0)&=[M J_1(\xi)+iN J_2(\xi)]e^{i\frac{kl}2(\cos\vartheta+\cos\vartheta_0)}\\
&\quad-g(3)[J_1(\xi)-iJ_2(\xi)]e^{-i\frac{kl}2(\cos\vartheta+\cos\vartheta_0)},
\end{aligned}\tag{15.10}\]
which are suitable for calculating the fringing field in the region $\varphi=-\frac\pi2$, $\frac\pi2<\vartheta;\vartheta_0<\pi$ $(\vartheta_0=\pi-\gamma)$. The quantities here equal
\[\xi=ka(\sin\vartheta+\sin\vartheta_0)\tag{15.11}\]
\[\left.\begin{aligned}
\left.\begin{matrix}\overline M\\M\end{matrix}\right\}&=\frac{\sin\frac\pi n}n\times\\
&\times\left(\mp\frac1{\cos\frac\pi n-\cos\frac{\pi-\vartheta-\vartheta_0}n}\pm\frac1{\cos\frac\pi n-\cos\frac{\vartheta+\vartheta_0}n}\right),\\
\left.\begin{matrix}\overline N\\N\end{matrix}\right\}&=\frac{\sin\frac\pi n}n\left(\frac2{\cos\frac\pi n-\cos\frac{\vartheta-\vartheta_0}n}\right.\\
&\left.\quad\mp\frac1{\cos\frac\pi n-\cos\frac{\pi-\vartheta-\vartheta_0}n}\mp\frac1{\cos\frac\pi n-\cos\frac{\vartheta+\vartheta_0}n}\right).
\end{aligned}\right\}\tag{15.12}\]
\end{scanpage}
